%-*-tex-*-
% Copyright Michael J. Ferguson, INRS-Telecommunications
% All rights reserved. 

%======== List macros ==========
% there is only one kind of list in this system, namely on that is
% right justified in the particular box that the list appears
% extensions may be made to include a center
% These list macros work relative to \rightskip and \leftskip defined
% outside the lists. In addition, they preserve the glue that already
% exists in \rightskip and \leftskip 

% the actual skip values are set in inrsdef.tex

\newskip\prlistskip %\prlistskip= 2ex plus 3pt minus 2pt
\newskip\prsublistskip %\prsublistskip = 1ex plus 2pt minus 1pt
\newskip\prsubsublistskip %\prsubsublistskip = .5ex plus 1pt minus 1pt

\newskip\polistskip %\polistskip= 2ex plus 1pt minus 1pt
\newskip\posublistskip %\posublistskip= 1ex plus 1pt minus 1pt
\newskip\posubsublistskip %\posubsublistskip= .5ex plus 1pt minus 1pt

\newskip\listitemskip %\listitemskip = 2ex plus 1pt minus 1pt
\newskip\sublistitemskip %\sublistitemskip =1ex plus .75pt minus .75pt
\newskip\subsublistitemskip %\subsublistitemskip = .5ex plus .5pt minus .5pt

\newskip\l@istitemskip 

\newdimen\listindent %\listindent = 3em
\newdimen\sublistindent %\sublistindent = 6em
\newdimen\subsublistindent %\subsublistindent = 9em
\newdimen\listitemmarksize  %\listitemmarksize = 1.5em

\newtoks\everylist
\newtoks\everysublist
\newtoks\everysubsublist
\newtoks\e@verylist
\newtoks\e@verysublist
\newtoks\e@verysubsublist

\newif\iff@irstli  % allows first li to set list parameters


% t@ightenlist will cause the list item spacing to drop to #1 of its
% original value 

%------ scales a skip .... makes plus glue .1 of skip, minus .1 ----
% #1 is the skip, #2 is the scale
\def\s@caleskip#1#2{\setbox0=\hbox{\hskip#1} 
                    \dimen0=.1\wd0 \dimen2=.1\wd0
                    #1 = #2\wd0 plus #2\dimen0 minus #2\dimen2}


\def\t@ightenlist#1{\s@caleskip{\listitemskip}{#1}
                    \s@caleskip{\sublistitemskip}{#1}
                    \s@caleskip{\subsublistitemskip}{#1}}
\def\tightenlist{\t@ightenlist{.5}}
\def\loosenlist{\t@ightenlist{2}} 

%sample list mark .... used for  creating lists where the listitem 
%mark size may be big or varied.
% This is one of the most useful features of this list package. 

\def\samplemark#1{\setbox0 = \hbox{#1\enspace}
                  \advance \listindent by -\listitemmarksize
                  \advance \sublistindent by -\listitemmarksize
                  \advance \subsublistindent by -\listitemmarksize
                  \listitemmarksize=\wd0
                  \advance \listindent by \listitemmarksize
                  \advance \sublistindent by \listitemmarksize
                  \advance \subsublistindent by \listitemmarksize}

%beginlist takes arbitrary parameters that affect the whole list such
%  as font changes ... they are local to the list and will be undone by
% the endlist
\newcount\l@iststate \l@iststate=0  % 0 -- outside lists , 1 in list, 2 sublist ...

%#1 - liststate, #2 - break value, #3 - listskip 
\newif\ifo@utsidelist \o@utsidelisttrue 
\newdimen\leftlistoffset 
\newdimen\rightlistoffset 
\def\b@eginlist #1#2#3{\par\penaltybreak{#2}{#3}\begingroup\parskip=0pt
                         \normalbaselines
                         \f@irstlitrue
                         \ifo@utsidelist \rightlistoffset =\rightskip
                            \leftlistoffset=\leftskip \o@utsidelistfalse\fi 
                          \the\b@@el\the\b@el\l@iststate=#1}

%\def\endlistmessage#1{<< endlist mismatch --- expected \the\l@iststate : 
%   actual #1 >> } see english/french titles ... #2 .. skip #3= penalty
\def\e@ndlist #1#2#3{\par \removelastskip 
                   \ifnum \l@iststate = #1 \else \message{\endlistmessage{#1}} \fi
                   \endgroup \penaltybreak{#3}{#2}}
% The first list item actually sets the list parameters. This allows for
% changes in at the beginning of a list in listitem indent and spacing.
\def\s@etlistparms{\iff@irstli \n@ewlistparms 
                          \advance\l@istindent by \leftlistoffset
                          \a@dvance{\leftskip}{\l@istindent}\f@irstlifalse
                      \fi}
\def\n@ewlistparms{\ifcase\l@iststate 
                    \or \let\l@istindent =\listindent
                        \let\l@istitemskip = \listitemskip
                    \or \let\l@istindent =\sublistindent
                        \let\l@istitemskip = \sublistitemskip
                    \or \let\l@istindent =\subsublistindent
                        \let\l@istitemskip = \subsublistitemskip \fi}

% #1 = itemskip #2 = leftindent #3= item flag
\def\listitem#1 {\par
                 \penaltybreak{-50}{\l@istitemskip}\s@etlistparms
                 \setbox0=\hbox{#1\enspace\hss}
                 \ifdim\wd0 < \listitemmarksize  \wd0= \listitemmarksize \fi
                 \noindent\hskip -\listitemmarksize
                 \box0\relax}
\def\beginlist {\let\b@@el=\e@verylist\let\b@el=\everylist 
                  \b@eginlist{1}{0}{\prlistskip}}
\def\beginsublist {\let\b@@el=\e@verysublist\let\b@el=\everysublist 
                  \b@eginlist{2}{100}{\prsublistskip}}
\def\beginsubsublist {\let\b@@el=\e@verysubsublist\let\b@el=\everysubsublist
                    \b@eginlist{3}{200}{\prsubsublistskip}}
                   

\def\endlist {\e@ndlist{1}{\polistskip}{-200}}
\def\endsublist {\e@ndlist{2}{\posublistskip}{-100}}
\def\endsubsublist {\e@ndlist{3}{\posubsublistskip}{-50}}

% ----- These are official short forms for the various items ----- 

\let\li =\listitem
\let\bl =\beginlist
\let\bsl=\beginsublist
\let\bssl=\beginsubsublist
\let\el=\endlist
\let\esl = \endsublist
\let\essl = \endsubsublist
